Los datos se han convertido en un activo valioso para las organizaciones que continuamente buscan maneras de mejorar la toma de decisiones y la eficiencia operativa. Sin embargo, a menudo se subutilizan en los procesos comerciales, donde los equipos priorizan el contacto directo con \textit{leads} sobre el análisis de su comportamiento a lo largo del ciclo de ventas. El presente proyecto propone un cambio hacia un enfoque basado en analítica de datos, mediante el cual los equipos comerciales pueden utilizar la información disponible para comprender mejor los patrones de comportamiento de los clientes y orientar sus decisiones. En el contexto de las empresas desarrolladoras inmobiliarias en la Ciudad de Guatemala, los proyectos residenciales generan un alto volumen de contactos, lo que implica para los asesores comerciales el reto de gestionar el tiempo y los recursos disponibles entre múltiples clientes potenciales. Tratar todos los \textit{leads} de la misma manera no necesariamente conduce a los mismos resultados, especialmente en el sector inmobiliario, donde los ciclos de compra suelen ser largos e involucran múltiples interacciones previas a la conversión. Sin una guía analítica para establecer un orden de atención, los esfuerzos comerciales pueden resultar ineficientes.

El objetivo del proyecto fue desarrollar y validar modelos predictivos aplicados a la gestión comercial de proyectos residenciales de vivienda en la Ciudad de Guatemala. El enfoque principal se centró en actividades de mercadeo y gestión comercial de vivienda, para las cuales se desarrolló un modelo de priorización de \textit{leads} basado en la probabilidad de conversión, complementado con un análisis de segmentación mediante \textit{k-means} para identificar perfiles de clientes. Este enfoque proporciona un criterio basado en datos para dirigir los recursos disponibles hacia las oportunidades con mayor potencial. Más allá del análisis predictivo, el proyecto incluyó el mapeo del proceso comercial mediante notación BPMN y la identificación de la etapa en la que debía incorporarse el sistema desarrollado. El análisis empírico se realizó con información de seis proyectos residenciales específicos, utilizados como caso de aplicación.

Los resultados obtenidos respaldan la viabilidad del enfoque propuesto como mecanismo para ordenar la cartera comercial. El modelo final logra recuperar el 70\% de las conversiones reales al concentrar la gestión en menos de un tercio de la cartera y, en el primer decil del ordenamiento, reúne cuatro veces más conversiones que una selección realizada al azar. Su principal aporte, por tanto, no radica en determinar de forma individual qué prospectos se convertirán, sino en establecer una secuencia de atención que concentre los esfuerzos en aquellos con mayor probabilidad de hacerlo. En este contexto, la precisión obtenida en el punto de operación fue del 7.9\%, lo que refleja también la baja frecuencia de conversiones en los datos analizados. Asimismo, las características relacionadas con el origen del \textit{lead} aportaron mayor capacidad predictiva que las variables demográficas, lo que evidencia la relevancia del contexto de captación para ordenar la cartera.

El documento se organiza de la siguiente manera: los capítulos iniciales sitúan el trabajo en relación con la investigación previa, establecen los objetivos, el alcance y las limitaciones, exponen la justificación y desarrollan el marco conceptual que sustenta el análisis. El capítulo de metodología describe la aplicación de las seis fases del marco CRISP-DM. A continuación, el capítulo de resultados presenta los hallazgos del análisis exploratorio, la segmentación, el modelo predictivo y su evaluación operativa. El capítulo de discusión los interpreta y examina sus implicaciones operativas y las limitaciones del estudio. Finalmente, se presentan las conclusiones y recomendaciones, mientras que los anexos documentan la fórmula completa del modelo y la trazabilidad de las cifras reportadas.